\chapter{External dependencies}
As mentioned in chapter \ref{planningTheProject}, CTaps is reliant on several
external libraries, both for testing and for core functionality. The largest
and most essential of these are libuv \cite{Libuv2025}. Libuv provides the central
event loop which allows CTaps to be asynchronous. It also provides several quality
of life features, such as built in support for asynchronous interactions with TCP
and UDP, through the \texttt{uv\_tcp\_t} and \texttt{uv\_udp\_t} handles respectively.
When adding protocols which libuv does not have native support for, such as QUIC,
it offers a custom work handle \texttt{uv\_work\_t}, which can be used to define
a custom piece of work. This can be used to define both a long-lived QUIC-based
server, or a short interaction with an existing server. Which underlying handle
is actually used in libuv is naturally abstracted away in CTaps, and for the
end-user of the library the interaction is the same no matter what. \todo {very informal lol}
